% Tipo de documento. En este caso es un art�culo, para folios A4, tama�o de la fuente 11pt y con p�gina separada para el t�tulo
\documentclass[a4paper,11pt,titlepage]{article}

% Carga de paquetes necesarios. OrdenesArticle es un paquete personalizado
\usepackage[spanish]{babel} 
\RequirePackage[T1]{fontenc}
\RequirePackage[ansinew]{inputenx} 
\usepackage[spanish,cap,cont,title,fancy]{OrdenesArticle}
\usepackage{lmodern}
\usepackage{graphicx}
\usepackage{hyperref}
\usepackage[usenames,dvipsnames]{color}
\usepackage{color}
\usepackage{ifthen}
\usepackage{longtable}
\hypersetup{bookmarksopen,bookmarksopenlevel=4,linktocpage,colorlinks,urlcolor=blue,citecolor=blue,
						linkcolor=blue,filecolor=blue,pdfnewwindow,
						pdftitle={An�lisis de un sistema colaborativo: Skype},
						pdfauthor={Juan Andrada Romero, Jose Domingo L�pez L�pez},
						pdfsubject={Sistemas para la Colaboraci�n}}


% Macro para definir una lista personalizada 
\newenvironment{milista}%
{\begin{list}{\textbullet}%
{\settowidth{\labelwidth}{\textbullet} \setlength{\leftmargin}{\dimexpr\labelsep+\labelwidth+5pt}
\setlength{\itemsep}{\dimexpr 0.5ex plus 0.25ex minus 0.25ex}
\setlength{\parsep}{\itemsep}
\setlength{\partopsep}{\itemsep}
\addtolength{\topsep}{-7.5pt}
}}%
{\end{list}}

% Macro para insertar una imagen
%       Uso: \imagen{nombreFichero}{Factor escala}{Caption (leyenda)}{Label (identificador para referenciarla)}
% -------------------------------------------------------------------------------------------------------------
\def\imagen#1#2#3#4{
 \begin{figure}[h]
 \begin{center}
   \scalebox{#2}{\includegraphics{#1}}
 \caption {#3}
 \label{#4}
 \end{center}
 \end{figure}
}

\begin{document}

% En las p�ginas de portada e �ndices, no hay encabezado ni pie de p�gina
\pagestyle{empty} 

% Se incluye la portada
\begin{titlepage}
	\begin{center}
  	{\LARGE UNIVERSIDAD DE CASTILLA-LA MANCHA} \\
  	\bigskip
  	{\Large ESCUELA SUPERIOR DE INFORM�TICA} \\
  	\vspace{20mm}
  	\includegraphics[scale=0.45, keepaspectratio]{./images/esi_bw.png} \\
  	\vspace{20mm}
  	{\Huge \textbf{Almacenamiento y Recuperaci�n de la Informaci�n}}\\
  	\vspace{10mm}
  	{\LARGE \textsc{\textbf{M�dulo de importaci�n e indexaci�n de documentos}}}\\
  	\vspace{30mm}
  	\today
  	\vspace{40mm}
  	\begin{flushleft}
  		{\large Juan Andrada Romero}\\
  		\vspace{1mm}
  		{\large Jose Domingo L�pez L�pez}\\		
  	\end{flushleft}
	\end{center}
\end{titlepage}
\clearpage

% Texto del reverso de la portada
%%\mbox{}
%%\vspace{18cm}
%%\begin{small}
% Se ajusta la separaci�n entre p�rrafos
%%\parskip=10pt 

%%\copyright~ 2008/2009 Juan Andrada Romero. Universidad de Castilla La Mancha, Escuela Superior de Inform�tica de Ciudad Real.

%%Se permite la modificaci�n, copia y distribuci�n de este documento, seg�n la licencia de documentaci�n GNU (\url{http://www.gnu.org}).

%%Este documento fue compuesto con \LaTeX{}. Im�genes generadas con Power Point y Gimp.
%%\end{small}

%%\newpage

% En las p�ginas de �ndices y prefacio, se utiliza numeraci�n romana
\pagenumbering{Roman}

% Se crea el �ndice
\tableofcontents
% Se pasa p�gina y se a�ade el �ndice de figuras al �ndice principal
%\clearpage\phantomsection
%\addcontentsline{toc}{section}{\listfigurename}
% Se crea el �ndice de figuras
%\listoffigures
% Si se quiere crear un �ndice de tablas se pondr�a: \listoftables

\newpage

% Se ajusta la separaci�n entre p�rrafos
\parskip=10pt

% Se a�ade el prefacio al �ndice
%\clearpage\phantomsection
%\addcontentsline{toc}{section}{Prefacio}

% Comienza el contenido del documento. Se utilizan n�meros ar�bigos y el encabezado y pie de p�gina personalizado
\pagenumbering{arabic}
\pagestyle{fancy}

\section{Introducci�n}


\section{Clasificaci�n}

\subsection{Dimensiones espacio-temporal}
Atendiendo a la clasificaci�n en tiempo y espacio propuesta por Johansen, \textit{Skype} puede clasificarse de la siguiente forma:

%\begin{milista}
%	\item \textbf{Mismo lugar, diferente tiempo}: aunque el prop�sito principal de \textit{Skype} es la comunicaci�n por voz en tiempo real utilizando el protocolo VoIP (voz por IP), se puede clasificar en esta categor�a porque tambi�n permite que los usuarios utilicen el chat que dicha herramienta incorpora, recibiendo los mensajes a posteriori si el usuario se encuentra desconectado. Es por ello que se trata, en esete caso particular, de una \textbf{interacci�n as�ncrona}. ---- Realmente el chat es aqui o en \textbf{Diferente lugar, diferente tiempo}?? ------
%	\item \textbf{Diferente lugar, mismo tiempo}: \textit{Skype} se puede clasificar dentro de esta categor�a, ya que las llamadas de voz en tiempo real son s�ncronas, pero cada usuario puede encontrarse en diferentes lugares. Del mismo modo ocurre con el chat, cuando los usuarios participantes est�n conectados, ya que los mensajes enviados se env�an a todos ellos.
%\end{milista}


% CHAT: diferente lugar, mismo tiempo. Interacci�n distribuida sincrona. Usuarios conectados (mismo tiempo) y diferente lugar (cada uno est� en su casa)
% TELEFONO: diferente lugar, mismo tiempo. Interacci�n distribuida sincrona. Usuarios conectados (mismo tiempo) y diferente lugar (cada uno est� en su casa)
% Videoconferencia: diferente lugar, mismo tiempo. Interacci�n distribuida sincrona. Usuarios conectados (mismo tiempo) y diferente lugar (cada uno est� en su casa)

Por otra parte, atendiendo a la clasificaci�n espacio-temporal propuesta por Grudin, la herramienta se puede clasificar en las siguientes categor�as:
%
%\begin{milista}
%	\item \textbf{Mismo lugar, diferente tiempo (impredecible)}: ---- Como dice pizarra, ��entraria aqui el chat, como en el caso anterior?? -----
%	\item \textbf{Diferente lugar(impredecible), mismo tiempo}: --- Videoconferencia y llamadas
%\end{milista}

% Ambiguedad porq existe para dispositivos moviles, por lo que no hay una forma de clasificar si la aplicacion es predecible o impredecible. 
% Diferente lugar, mismo tiempo.

\subsection{Dominio de aplicaci�n}
%ellis

A: Es un sistema de mensaje porq se envian mensaje de texto (chat), sonido (llamadas) y video (videoconferencia)
D: Es un sistema de conferencias. Es del tipo D4: Conferencia de escritorio.






	
	
\section{Descripci�n de los requisitos funcionales}

\subsection{Participaci�n}

% Sesion que se refiere a la sesion del servidor

1.1. Predefinido porq hay q estar registrado e inicar sesion en el servidor
1.2. Permitir entradas posteriores en el contexto del servidor
Introducir usuarios al inicio de la sesi�n, en el contexto de crear grupos de chats, llamadas y videoconferencias. Tambien es permitir entradas posteriores.
1.3. Permitir salir antes.
1.4. Aceptada, porq el usuario tiene que aceptar la invitaci�n que envia uno de los usuarios del grupo.
1.5. Segun el contexto de uso de la aplicaci�n. Para un uso dom�stico, normalmente es de 2 participantes y, ocasionalmente, peque�o, pero para un uso profesional y empresarial, suele ser peque�o grande.
1.7. Conjunta. Ejemplo de chat colaborativo
1.8. Textual
Voz y video
1.10. Persistente el log de chats y de acciones. Tb los contactos
No persistente audio y video
1.11 Depende de la tarea (chat, llamada...)
1.12 Id�ntico en el chat, log y llamada
Diferente en la videoconferencia (pantallita de cara)

\subsection{Interacci�n}

2.1. Sincrono
2.2. Distribuida
2.3. Ambos, dependiendo del contexto de uso de la aplicaci�n (empresarial, domestico)
2.4. Ambos, dependiendo del contexto de uso de la aplicaci�n (empresarial, domestico)
2.5. Democr�tico, porq todos los usuarios pueden realizar las mismas acciones (ejemplo IRC:administradores, etc)
2.6. Ambos, dependiendo del contexto de uso de la aplicaci�n (empresarial, domestico)
2.7. Estructurada ?????
2.8. Identificada
2.9. Simultanea

\subsection{Coordinaci�n}

% No puede haber mucha cosas porq no es una aplicacion de coordinacion

\subsection{Distribuci�n y soporte}

4.1. Centralizada en el servidor (logs, perfiles, contactos, compras, credito, etc.)
4.2. Hibrida, porq las peticiones y el login pasn por el servidor pero al establecer ya la comunicaci�n, el procesamiento se realiza en los clientes (por ejemplo, el video)


\subsection{Notificaci�n de eventos}

5.1. Consciente (de la interaccion)
5.3. Lista de usuarios con iconos de estado

\subsection{Visualizaci�n}

6.1. WYSIWIS relajada (chat)
WYSIWIS relajada (llamadas, porq la duracion es la misma pero la configuracion de la llamada no)
No aplicable para videoconferencias
6.2. Privada (porq cada cliente ejecuta su aplicacion)
6.3. Ninguna
6.4. ????
6.5. Si, para el chat
6.6. Actualizar el estado de los usuarios, avatar, log de eventos

\section{Explicaci�n del funcionamiento} \label{manual}

% Pantalla compartida a traves de una videoconferencia
    % Pantalla compelta o sleccion (una parte de la pantalla)
% Compartir archivos 
% Informacion de la calidad de las llamadas

% Comentar otras funcionalidades no colaborativas, como compra de credito, gestion del directorio y extrass
\section{Cr�ticas y propuestas de posibles mejoras}

% Secci�n de ap�ndices
% Se a�ade el titulo de "Ap�ndices" al �ndice
%\addcontentsline{toc}{section}{Ap�ndices}

%\input{Apendices.tex}

% A�adimos la bibliografia al �ndice
\clearpage\phantomsection
\addcontentsline{toc}{section}{Referencias}

% Bibliograf�a. En este caso se usa BibTeX
%\bibliographystyle{plain}
%\pagestyle{plain} 
%\bibliography{Bibliografia}

\end{document}
